% ~10 min talk: conceptual changes since forking the Juliano 3D-PBL branch,
% restricted to the full closure (pbl3d_opt=2). Build: tectonic talk_3dpbl_changes.tex
\documentclass[10pt,aspectratio=169]{beamer}
\usetheme{default}
\usecolortheme{seahorse}
\setbeamertemplate{navigation symbols}{}
\setbeamertemplate{footline}[frame number]
\setbeamerfont{frametitle}{series=\bfseries,size=\large}
\setbeamertemplate{itemize items}[circle]
\usepackage{amsmath,booktabs}
\newcommand{\avg}[1]{\langle #1 \rangle}
\newcommand{\pd}[2]{\frac{\partial #1}{\partial #2}}
\graphicspath{{/gpfs/data/fs72996/ewahl/plot_output/diagnostics/}}

\title{A full 3D Mellor--Yamada closure over real Alpine terrain}
\subtitle{What it took to make the forked scheme run --- and improve}
\author{Elias Wahl}
\institute{University of Innsbruck}
\date{September 2026}

\begin{document}

\begin{frame}
  \titlepage
\end{frame}

% ---------------------------------------------------------------- intro 1
\begin{frame}{Why a 3D closure at all}
  A 1D PBL scheme assumes each column is horizontally homogeneous. Of the nine
  second moments in the mean equations
  \[
    \pd{U_i}{t} + U_j \pd{U_i}{x_j}
      = \dots - \pd{\avg{u_i u_j}}{x_j},
    \qquad
    \pd{\Theta}{t} + U_j \pd{\Theta}{x_j} = -\pd{\avg{u_j\theta}}{x_j},
  \]
  only $\avg{uw},\ \avg{vw},\ \avg{w\theta}$ survive that assumption --- three
  numbers per level, acting through their vertical derivative alone.
  \vspace{1ex}

  At $\Delta x = 500$~m in the Inn Valley this fails twice:
  \begin{itemize}
    \item \textbf{Grey zone}: $\Delta x \sim z_i$ --- the energetic eddies are
      neither subgrid nor resolved; horizontal gradients of the turbulence
      statistics are as large as vertical ones.
    \item \textbf{Slopes of $30^\circ$+}: a katabatic jet produces shear that
      projects onto \emph{every} strain component; a scheme that only sees
      $\partial U/\partial z$ is structurally blind to most of its production.
  \end{itemize}
  \vspace{1ex}
  The 3D PBL scheme (Kosovi\'c et al.\ 2020, Juliano et al.\ 2022) parameterizes
  all six stresses and all three components of each scalar flux, and applies
  their full 3D divergence --- replacing both the 1D PBL scheme \emph{and} the
  horizontal diffusion. There is no eddy-viscosity backstop behind it.
\end{frame}

% ---------------------------------------------------------------- intro 2
\begin{frame}{The full closure in one slide (\texttt{pbl3d\_opt}=2)}
  Mellor--Yamada level 2.5: one prognostic variable,
  $q^2 = \avg{u^2}+\avg{v^2}+\avg{w^2}$ (twice the TKE, m$^2$\,s$^{-2}$),
  plus a master length scale $l$ (m). All other moments are \emph{algebraic}.
  \vspace{1ex}

  The full variant keeps all nine velocity gradients and all three of
  $\nabla\Theta_v$: the moment system no longer decouples, and at every grid
  point, every substep, a dense $10\times10$ linear system is solved
  \[
    \mathsf{A}\,\mathbf{x} = \mathbf{b},
    \qquad
    \mathbf{x} = \big(\avg{u^2},\avg{v^2},\avg{w^2},\ \avg{uv},\avg{uw},\avg{vw},\
    \avg{u\theta_v},\avg{v\theta_v},\avg{w\theta_v},\ \avg{\theta_v^2}\big)^{\!\top}
  \]
  with $\mathsf{A} = \tau^{-1}\mathsf{I} + \mathsf{S}$: return-to-isotropy rates
  $\propto q/l$ on the diagonal, mean strain off it, buoyancy $\beta g$ coupling
  the blocks. A $4\times4$ moisture block follows.
  \vspace{1ex}
  \begin{itemize}
    \item Scalar-flux rows have \emph{no source of their own} --- fluxes exist
      only because stresses act on scalar gradients.
    \item Near the turbulence cutoff the diagonal shrinks with $q$ while strain
      does not: strained, weakly turbulent air drives the system toward
      singularity --- this is where the fork stood.
  \end{itemize}
\end{frame}

% ---------------------------------------------------------------- change 1
\begin{frame}{Change 1 --- an answer at every grid point: the realizability stack}
  \textbf{At the fork} the full closure did not survive real terrain: when the
  LAPACK solve failed, the solution vector was \emph{uninitialized memory}; when
  it succeeded, nothing checked that the returned stress tensor was physically
  possible. Unrealizable solutions are not rare --- 6--40\,\% of points; singular
  systems only $\sim$0.005\,\%.
  \vspace{1ex}

  \textbf{Concept: never hand the dynamics an impossible tensor.} Four tiers,
  each an \emph{exact} solution of the closure, at shorter $l$ if needed:
  \begin{enumerate}
    \item equilibrated double-precision solve, condition estimate checked;
    \item Durbin strain limit $l \le (2\,c_{\mathrm{lim}}/b_1)\,q/|S|$ ---
      keeps the closure inside its own validity range, inactive at equilibrium;
    \item while the solution is degenerate or the stress tensor is not positive
      semi-definite: halve $l$, resolve (full Sylvester test as acceptance);
    \item exact realizability completion: variance floor, trace identity,
      Cauchy--Schwarz, determinant-based covariance shrink.
  \end{enumerate}
  \vspace{1ex}
  Verified on 200\,000 adversarial random states: zero disagreements with a true
  positive-semi-definite test, zero bound violations after completion.
\end{frame}

% ---------------------------------------------------------------- change 2
\begin{frame}{Change 2 --- the length scale must come from the turbulence}
  Three defects, one theme: \textbf{$l$ was set by things that are not the flow.}
  \vspace{1ex}
  \begin{itemize}
    \item \textbf{The asymptotic scale $l_0$ was set by the model lid.} MY74's
      $l_0 \propto \int q z\,dz / \int q\,dz$ converges because $q\to0$ aloft ---
      but $q^2$ is floored, so every quiescent level contributes and
      $l_0 \to \alpha H/2$: 250\,m for a 5\,km top, 1500\,m for 30\,km.
      Fixed by excluding floor-level air from the integral.
    \item \textbf{No buoyancy limit.} Above the boundary layer
      $\tau = l/q$ reached $10^4$\,s, $N\tau \sim 100$ --- eddies a hundred
      buoyancy periods long. The Deardorff limit $l \le c\,q/N$ now caps the
      master scale in stable air.
    \item \textbf{One master scale everywhere.} Mixing length and dissipation
      length had diverged; recoupling them removed a nocturnal runaway in which
      $q^2$ grew without bound over slopes. The strain cap of tier 2 is
      load-bearing there: relaxing it brings the runaway back within 45 min.
  \end{itemize}
  \vspace{1ex}
  $N$: Brunt--V\"ais\"al\"a frequency (s$^{-1}$); $\tau$: eddy turnover time.
\end{frame}

% ---------------------------------------------------------------- change 3
\begin{frame}{Change 3 --- subgrid energy must be paid by the resolved flow}
  \textbf{Concept: the closure's $q^2$ production and the kinetic energy the
  resolved flow loses are the same physical transaction} --- so audit them
  against each other (WRFlux budgets made this measurable).
  \vspace{1ex}

  The audit found a leak: on slopes, the horizontally-paired shear production
  --- about a third of \emph{all} production, essentially all of it on slopes ---
  was credited to $q^2$ while the corresponding momentum tendency was tapered
  away. $\sim$90\,\% of that energy was never extracted from the resolved wind.
  \vspace{1ex}
  \begin{center}
  \begin{tabular}{lcc}
    \toprule
     & before pairing fix & after \\
    \midrule
    energy-closure residual & $+14$ to $+37\,\%$ & $+0.3\,\%$ \\
    nocturnal $q^2$ (fraction of MYNN, 04\,UT) & 0.33 & 0.17 \\
    slope-bin structure of $q^2$ & strong & flat \\
    \bottomrule
  \end{tabular}
  \end{center}
  \vspace{1ex}
  The slope-dependent $q^2$ structure \emph{was} the spurious source. Free
  energy from an accounting mismatch looks exactly like turbulence physics
  until the books are balanced.
\end{frame}

% ---------------------------------------------------------------- change 4
\begin{frame}{Change 4 --- separate the closure from the host model}
  Two upstream WRF bugs spent weeks disguised as closure physics:
  \vspace{1ex}
  \begin{itemize}
    \item \textbf{Surface-layer lookup table} out-of-bounds under strong
      stability (crash) --- patched at the table bound, not papered over.
    \item \textbf{Undefined albedo ($-9999$) handed to radiation}: with
      topographic shading on, the Noah-MP driver gives shortwave radiation its
      fill value wherever the surface gets no direct beam --- RRTMG then
      \emph{cools} shaded ground at $-80$\,K/h. The morning fog, the $-11$\,K
      cold bias, the drainage jets and the 07:54 crash were all downstream of
      one number --- and were withdrawn as closure results.
  \end{itemize}
  \vspace{1ex}
  \textbf{Method that came out of it:} when a budget term is \emph{impossible}
  (shortwave cooling), stop hypothesizing physics and read the raw model arrays
  at the offending cells --- found in an hour from a restart file, after a day
  of closure hypotheses.
\end{frame}

% ---------------------------------------------------------------- change 5
\begin{frame}{Change 5 --- judging a 3D closure in the grey zone}
  \begin{columns}[c]
    \column{0.52\textwidth}
    The 3D run's subgrid $q^2$ is 0.28 of MYNN's in the morning mixed layer.
    \textbf{That is a partition, not a deficit}: the closure leaves to the
    resolved flow what the resolved flow can do.
    \vspace{1ex}
    \begin{itemize}
      \item subgrid $+$ resolved kinetic energy: 3.9 vs.\ MYNN's 3.7
        m$^2$\,s$^{-2}$ at 43\,m;
      \item the 3D run resolves 85--91\,\% of it (MYNN: $\le$20\,\%);
      \item its mixed layer is deeper, its total entrainment flux larger.
    \end{itemize}
    \vspace{1ex}
    \textbf{Rule:} in unstable air, compare subgrid \emph{and} resolved
    energy, never subgrid alone. (In stable air there is no resolved
    turbulence at 500\,m --- there the subgrid number stands alone.)
    \column{0.46\textwidth}
    \includegraphics[width=\textwidth]{gz_partition/gz_partition_18.png}
    {\footnotesize Entrainment heat flux at the CBL top, valley floor: the 3D
    closure (purple) carries it resolved; MYNN (green) subgrid.}
  \end{columns}
\end{frame}

% ---------------------------------------------------------------- change 6
\begin{frame}{Change 6 --- scalar realizability opens the daytime valley}
  \begin{columns}[c]
    \column{0.46\textwidth}
    The tier-2 acceptance tested the \emph{momentum} tensor only. Rejecting a
    state for its stresses also discarded its heat flux --- and on the heated
    valley walls that killed the upslope heat transport: the boundary layer
    stopped growing at $\sim$300\,m.
    \vspace{1ex}

    Accepting the scalar block on its own realizability
    (\texttt{pbl3d\_t2\_scalar}):
    \vspace{0.5ex}
    \begin{center}
    \begin{tabular}{lc}
      \toprule
      $h_\theta$, Kolsass 11\,UT & \\
      \midrule
      observed (radiosonde) & 955\,m \\
      closed wall & 317\,m \\
      scalar acceptance & 746\,m (11\,UT) \\
       & 1136\,m (12\,UT) \\
      \bottomrule
    \end{tabular}
    \end{center}
    \vspace{0.5ex}
    Validated in every regime; bit-for-bit identical when off.
    \column{0.52\textwidth}
    \includegraphics[width=\textwidth]{t2s_sounding_11utc.png}
    {\footnotesize Valley $\theta$ profiles $\sim$11\,UT: the closed wall (red)
    never mixes the valley; the fix (blue) reaches the observed structure with
    $\le$1\,h lag.}
  \end{columns}
\end{frame}

% ---------------------------------------------------------------- stable regime
\begin{frame}{Where the science stands --- the stable regime is the frontier}
  \begin{columns}[c]
    \column{0.50\textwidth}
    \includegraphics[width=\textwidth]{x10_fig4_q2_timeheight.png}
    {\footnotesize Median subgrid $q^2$, night into morning: MYNN mixes the
    whole valley volume all night; the 3D closure carries an order of magnitude
    less --- its own stable-regime equilibrium, flat above Ri~$\approx0.3$.}
    \column{0.48\textwidth}
    \begin{itemize}
      \item The nocturnal $q^2$ deficit against MYNN is the closure's stable
        equilibrium --- not spin-up, not initialization. And MYNN is the
        \emph{control, not the truth}: it over-mixes the night and
        under-forecasts the valley wind.
      \item No constant is tuned before a full diurnal cycle is held against
        observations --- every new switch defaults to the previous behaviour,
        bit for bit.
    \end{itemize}
  \end{columns}
\end{frame}

% ---------------------------------------------------------------- cold pool
\begin{frame}{The cold pool: a model equilibrium, not a closure property}
  \begin{columns}[c]
    \column{0.42\textwidth}
    \includegraphics[width=\textwidth]{x10_fig1_theta_kolsass_05.png}
    {\footnotesize Kolsass 05\,UT: every configuration converges to the same
    too-warm pool.}
    \column{0.56\textwidth}
    The $\sim$3.5\,K warm bias of the mature valley cold pool is
    \textbf{invariant} --- measured, not assumed --- to:
    \begin{itemize}
      \item the \emph{closure} (MYNN twin: same deficit),
      \item the \emph{initial state} (evening start vs.\ 01\,UT start),
      \item the \emph{forcing resolution} (11 vs.\ 65 ICON levels: a 1.4\,K
        warmer evening start is eaten by 22\,UT; identical pool at 01\,UT).
    \end{itemize}
    \vspace{1ex}
    The measured levers sit elsewhere:
    \begin{itemize}
      \item the 6th-order filter warms the pool floor across
        terrain-following surfaces ($+0.12$--$0.15$\,K/h $\times$ 10\,h night
        $\approx$ the sunrise deficit);
      \item missing evening cold-air production/import (19--21\,UT);
      \item surface coupling: skin held warm, exchange that refuses to die
        at moderate Ri.
    \end{itemize}
  \end{columns}
\end{frame}

% ---------------------------------------------------------------- forcing
\begin{frame}{A late lesson: the forcing is part of the model}
  The driving ICON data entered as an 11-pressure-level product. Two
  consequences, both found masquerading as model biases:
  \vspace{1ex}
  \begin{itemize}
    \item no level between $\sim$580 and $\sim$1520\,m ASL --- the nocturnal
      cold pool was \emph{interpolation-filled} ($+2$\,K) in every analysis-time
      initial state;
    \item the afternoon up-valley jet was smeared out of the initial condition
      entirely (2.3 vs.\ 5.6\,m/s observed at first contact).
  \end{itemize}
  \vspace{1ex}
  Rebuilt conversion on the 65 native ICON levels; a 36-level pressure ladder
  proved fully equivalent (twin runs agree to tenths). What remains of the
  daytime wind deficit ($\sim$60\,\% of observed) is \emph{in the forcing
  itself}, measured directly in the boundary files --- at 500\,m the
  along-valley momentum is largely imposed at the lateral boundaries.
\end{frame}

% ---------------------------------------------------------------- summary
\begin{frame}{Summary --- the conceptual arc}
  \begin{enumerate}
    \item \textbf{Realizability as the contract}: every grid point returns an
      exact, physically possible solution of the closure --- momentum
      \emph{and} scalar blocks judged on their own terms.
    \item \textbf{The length scale belongs to the turbulence} --- not the lid,
      not the floor value, not a second bookkeeping copy.
    \item \textbf{Energy accounting closes the loop}: subgrid production is
      real only if the resolved flow pays it.
    \item \textbf{Blame is established by measurement}: host-model bugs,
      forcing artifacts and closure physics separated by budgets and raw
      arrays, not by plausibility.
    \item \textbf{Grey-zone judgement}: subgrid + resolved, against
      observations --- with MYNN as control, not truth.
  \end{enumerate}
  \vspace{1ex}
  Discipline that made it possible: every switch default-off and bit-for-bit;
  statistical comparison only; same-turn records of every decision.
  \vspace{1ex}

  Open: the stable-regime equilibrium vs.\ observed nocturnal turbulence;
  the three cold-pool levers; ICON-terrain and forcing-resolution chains
  currently running.
\end{frame}

\end{document}
